\documentclass[10pt,journal,compsoc,twocolumn]{IEEEtran}
\usepackage[utf8]{inputenc}
\usepackage{amsmath,amssymb,amsfonts}
\usepackage{graphicx}
\usepackage{booktabs}
\usepackage{cite}
\usepackage{hyperref}
\usepackage{microtype}

\begin{document}

\title{Task C3: Multi-Reference Image Captioning via CNN-LSTM Fusion and Beam Search Optimization}

\author{
    \IEEEauthorblockN{Muhammad Abdullah}\\
    \IEEEauthorblockA{Department of Intelligent Systems, The University of Lahore, Pakistan\\
    Roll No: 70178208 \quad|\quad Section: BSRAI-5C\\
    GitHub: \url{https://github.com/MuhammadaAbdullah70178208/AI101-ML-Assignment-70178208}\\
    Instructor: Dr. Muhammad Awais, AI Architect}
}

\maketitle

\begin{abstract}
Cross-modal vision-to-language generation requires translating high-dimensional visual scenes into syntactically and semantically coherent natural language descriptions. This investigation addresses Task C3 of the AI-101 course syllabus, operationalizing the sequence generation framework and evaluation methods from Modules~01 and~02 on the multi-reference Flickr8k benchmark ($N = 2,000$ images, 5 references per image). We implement a CNN-LSTM merge architecture and benchmark decoding strategies against a pre-trained Vision Transformer baseline (ViT-GPT2). Our findings reveal that expanding the decoding search space from Greedy ($k=1$) to Beam Search ($k=5$) improves corpus BLEU-4 from 15.2\% to 20.8\% (+36.8\% relative improvement) while increasing vocabulary richness. Furthermore, the Vision Transformer model achieves 25.4\% BLEU-4 and 68.8\% BLEU-1, demonstrating that patch-based visual self-attention substantially reduces descriptive hallucinations in complex multi-object scenes.
\end{abstract}

\section{Introduction}
Automated image description bridges computer vision and natural language processing, enabling technologies such as assisted accessibility for the visually impaired, automated photo organization, and multi-modal content indexing. Unlike discrete image classification, captioning requires a model to perceive spatial relationships, ground entities, and generate sequential syntax adhering to target linguistic grammar. In this study, we implement an end-to-end captioning lifecycle, systematically analyzing decoding search width and architectural paradigms.

\section{Problem Formulation and Evaluation Protocol}
Given an input photographic image $\mathbf{I} \in \mathbb{R}^{H \times W \times C}$, our objective is to generate a descriptive sequence of word tokens:
\begin{equation}
\mathbf{y}^* = \arg\max_{\mathbf{y}} P(\mathbf{y} \mid \mathbf{I}) = \arg\max_{\mathbf{y}} \prod_{t=1}^T P(y_t \mid y_{<t}, \mathbf{I})
\end{equation}
Because natural scenes admit multiple valid human descriptions, evaluation utilizes multi-reference modified n-gram precision across five ground-truth captions ($ref_1, \dots, ref_5$):
\begin{equation}
p_n = \frac{\sum_{g \in \text{hyp}} \min\left(C(g, \text{hyp}), \max_{r \in \text{refs}} C(g, r)\right)}{\sum_{g \in \text{hyp}} C(g, \text{hyp})}
\end{equation}
combined with an exponential brevity penalty: $\text{BLEU-N} = \text{BP} \cdot \exp\left(\sum_{n=1}^N w_n \log p_n\right)$.

\section{Dataset Characteristics and Splitting}
We utilize the multi-reference Flickr8k benchmark (2,000 images, 10,000 reference captions) spanning outdoor play, sports, urban transit, and nature.

\begin{table}[htbp]
\centering
\caption{Flickr8k Multi-Reference Benchmark Characteristics}
\label{tab:c3_stats}
\begin{tabular}{ll}
\toprule
\textbf{Specification} & \textbf{Value} \\
\midrule
Total Photographic Images & 2,000 scenes \\
Reference Captions / Image & 5 independent human annotations \\
Target Modality & Vision-to-Language Sequence Generation \\
Vocabulary Size & $V = 3,850$ unique word types \\
Partition Protocol & Disjoint 70\% Train ($1,400$) / 15\% Val ($300$) \\
                   & / 15\% Test ($300$ images) \\
\bottomrule
\end{tabular}
\end{table}

\begin{figure}[htbp]
\centering
\includegraphics[width=\columnwidth]{figures/c3_bleu_beam_search_comparison.png}
\caption{Task C3 Diagnostics: (Left) BLEU-1 to BLEU-4 metrics across decoding search regimes; (Right) Cross-entropy loss convergence over 15 training epochs.}
\label{fig:bleu_loss}
\end{figure}

Splitting follows a disjoint 70/15/15 protocol based on image IDs to ensure zero visual leakage. Text tokenization is fitted strictly on training captions.

\section{Model Architectures and Decoding Regimes}
We investigate two core architectures and multiple inference search widths:
\begin{enumerate}
    \item \textbf{CNN-LSTM Merge Model:} A convolutional feature extractor maps image pixels to a 512-dimensional visual vector $\mathbf{v} \in \mathbb{R}^{512}$. At each recurrent step $t$, the decoder combines visual features with word embeddings:
    \begin{equation}
    h_t = \text{LSTM}\left([e(y_{t-1}); \mathbf{v}], h_{t-1}\right)
    \end{equation}
    followed by a dense softmax projection over the vocabulary.
    \item \textbf{Greedy Decoding ($k=1$):} Choosing the locally optimal token $\hat{y}_t = \arg\max P(y_t \mid y_{<t})$.
    \item \textbf{Beam Search Decoding ($k \in \{3, 5, 8\}$):} Maintaining $k$ candidate hypotheses at each time step, pruning paths based on cumulative log-probabilities $\sum_{\tau=1}^t \log P(y_\tau \mid y_{<\tau})$.
    \item \textbf{ViT-GPT2 Benchmark:} Vision Transformer encoder paired with auto-regressive GPT-2 decoder.
\end{enumerate}

\section{Results and Performance Analysis}
Quantitative multi-reference BLEU evaluation across 300 unseen test scenes is summarized in Table~\ref{tab:results}.

\begin{table}[htbp]
\centering
\caption{Quantitative Performance on Unseen Test Partition ($N=300$)}
\label{tab:results}
\begin{tabular}{lcccc}
\toprule
\textbf{Architecture / Search} & \textbf{BLEU-1} & \textbf{BLEU-2} & \textbf{BLEU-3} & \textbf{BLEU-4} \\
\midrule
CNN-LSTM (Greedy $k=1$)   & 58.4\% & 38.6\% & 24.5\% & 15.2\% \\
CNN-LSTM (Beam $k=3$)     & 61.2\% & 41.8\% & 27.9\% & 18.5\% \\
\textbf{CNN-LSTM (Beam $k=5$)} & \textbf{63.5\%} & \textbf{44.1\%} & \textbf{30.2\%} & \textbf{20.8\%} \\
CNN-LSTM (Beam $k=8$)     & 63.6\% & 44.2\% & 30.3\% & 21.0\% \\
\midrule
ViT-GPT2 Benchmark        & 68.8\% & 49.5\% & 35.8\% & 25.4\% \\
\bottomrule
\end{tabular}
\end{table}

\begin{figure}[htbp]
\centering
\includegraphics[width=\columnwidth]{figures/c3_loss_and_diversity.png}
\caption{Task C3 Optimization: Corpus BLEU-4 versus beam width ($k \in \{1, 3, 5, 8\}$) illustrating diminishing returns beyond $k=5$.}
\label{fig:beam_curve}
\end{figure}

\section{Discussion and Error Analysis}
1. \textbf{Search Width Optimization:} Greedy search ($k=1$) frequently suffers from irreversible sub-optimal choices: an early incorrect verb or article permanently derails subsequent noun phrase generation. Transitioning from Greedy to Beam Search with $k=5$ yields a massive **+5.6 percentage point gain in BLEU-4 (15.2\% to 20.8\%)**.
2. \textbf{Pareto Efficiency:} Increasing beam width from $k=5$ to $k=8$ yields a negligible gain of +0.2\% BLEU-4 while increasing inference compute by 60\%, identifying $k=5$ as the optimal deployment threshold.
3. \textbf{Vision Transformer Advantage:} ViT-GPT2 achieved 25.4\% BLEU-4, outperforming CNN-LSTM by +4.6 points. Spatial self-attention patches enable the language model to condition on discrete localized visual objects (e.g., ``tennis ball'', ``wetsuit'') rather than a single pooled image vector.

\section{Conclusion}
This study demonstrated that recurrent CNN-LSTM architectures achieve strong multi-reference captioning performance, with beam search decoding providing a 36.8\% relative boost in BLEU-4 over greedy heuristics.

\begin{thebibliography}{1}
\bibitem{hodosh2013framing}
M.~Hodosh, P.~Young, and J.~Hockenmaier, ``Framing image description as a ranking task: Data, models and evaluation metrics,'' \emph{JAIR}, vol.~47, pp.~853--899, 2013.
\bibitem{vinyals2015show}
O.~Vinyals, A.~Toshev, S.~Bengio, and D.~Erhan, ``Show and tell: A neural image caption generator,'' in \emph{Proc. CVPR}, 2015.
\bibitem{awais2026introai}
M.~Awais, \emph{AI-101: Introduction to Artificial Intelligence (Course Slides: Modules 01 \& 02)}, Dept. of Intelligent Systems, The University of Lahore, 2026.
\end{thebibliography}

\end{document}
